\Chapter{95}

\Verse{1}
{
    \zR{话说焙茗在门口和小丫头子说宝玉的玉有了，那小丫头急忙回来告诉宝玉。}
}
{
    \eR{[English source not present in the checked-in Bencraft/Joly files.]}
}
{
}

\Verse{2}
{
    \zR{众 人听了，都推着宝玉出去问他。}
}
{
    \eR{[English source not present in the checked-in Bencraft/Joly files.]}
}
{
}

\Verse{3}
{
    \zR{众人在廊下听着。}
}
{
    \eR{[English source not present in the checked-in Bencraft/Joly files.]}
}
{
}

\Verse{4}
{
    \zR{宝玉也觉放心，便走到门口，问 道：“你那里得了？}
}
{
    \eR{[English source not present in the checked-in Bencraft/Joly files.]}
}
{
}

\Verse{5}
{
    \zR{快拿来。”}
}
{
    \eR{[English source not present in the checked-in Bencraft/Joly files.]}
}
{
}

\Verse{6}
{
    \zR{焙茗道：“拿是拿不来的，还是托人做保去呢。”}
}
{
    \eR{[English source not present in the checked-in Bencraft/Joly files.]}
}
{
}

\Verse{7}
{
    \zR{宝玉道：“你快说是怎么得的，我好叫人取去。”}
}
{
    \eR{[English source not present in the checked-in Bencraft/Joly files.]}
}
{
}

\Verse{8}
{
    \zR{焙茗道：“我在外头，知道林爷 爷去测字，我就跟了去。}
}
{
    \eR{[English source not present in the checked-in Bencraft/Joly files.]}
}
{
}

\Verse{9}
{
    \zR{我听见说在当铺里找，我没等他说完，便跑到几个当铺里 去。}
}
{
    \eR{[English source not present in the checked-in Bencraft/Joly files.]}
}
{
}

\Verse{10}
{
    \zR{我比给他们瞧，有一家便说‘有’。}
}
{
    \eR{[English source not present in the checked-in Bencraft/Joly files.]}
}
{
}

\Verse{11}
{
    \zR{我说：‘给我罢。’}
}
{
    \eR{[English source not present in the checked-in Bencraft/Joly files.]}
}
{
}

\Verse{12}
{
    \zR{那铺子里要票子。}
}
{
    \eR{[English source not present in the checked-in Bencraft/Joly files.]}
}
{
}

\Verse{13}
{
    \zR{我 说：‘当多少钱？’}
}
{
    \eR{[English source not present in the checked-in Bencraft/Joly files.]}
}
{
}

\Verse{14}
{
    \zR{他说：‘三百钱的也有，五百钱的也有。}
}
{
    \eR{[English source not present in the checked-in Bencraft/Joly files.]}
}
{
}

\Verse{15}
{
    \zR{前儿有一个人拿这么 一块玉，当了三百钱去；今儿又有人也拿一块玉当了五百钱去。’”}
}
{
    \eR{[English source not present in the checked-in Bencraft/Joly files.]}
}
{
}

\Verse{16}
{
    \zR{宝玉不等说完， 便道：“你快拿三百五百钱去取了来，我们挑着看是不是。”}
}
{
    \eR{[English source not present in the checked-in Bencraft/Joly files.]}
}
{
}

\Verse{17}
{
    \zR{里头袭人便啐道：“二 爷不用理他。}
}
{
    \eR{[English source not present in the checked-in Bencraft/Joly files.]}
}
{
}

\Verse{18}
{
    \zR{我小时候儿听见我哥哥常说，有些人卖那些小玉儿，没钱用便去当， 想来是家家当铺里有的。”}
}
{
    \eR{[English source not present in the checked-in Bencraft/Joly files.]}
}
{
}

\Verse{19}
{
    \zR{众人正在听得诧异，被袭人一说，想了一想，倒大家笑 起来，说：“快叫二爷进来罢，不用理那糊涂东西了。}
}
{
    \eR{[English source not present in the checked-in Bencraft/Joly files.]}
}
{
}

\Verse{20}
{
    \zR{他说的那些玉，想来不是正 经东西。”}
}
{
    \eR{[English source not present in the checked-in Bencraft/Joly files.]}
}
{
}

\Verse{21}
{
    \zR{宝玉正笑着，只见岫烟来了。}
}
{
    \eR{[English source not present in the checked-in Bencraft/Joly files.]}
}
{
}

\Verse{22}
{
    \zR{原来岫烟走到栊翠庵，见了妙玉，不及闲话，便 求妙玉扶乩。}
}
{
    \eR{[English source not present in the checked-in Bencraft/Joly files.]}
}
{
}

\Verse{23}
{
    \zR{妙玉冷笑几声，说道：“我与姑娘来往，为的是姑娘不是势利场中的 人。}
}
{
    \eR{[English source not present in the checked-in Bencraft/Joly files.]}
}
{
}

\Verse{24}
{
    \zR{今日怎么听了那里的谣言，过来缠我？}
}
{
    \eR{[English source not present in the checked-in Bencraft/Joly files.]}
}
{
}

\Verse{25}
{
    \zR{况且我并不晓得什么叫‘扶乩’。”}
}
{
    \eR{[English source not present in the checked-in Bencraft/Joly files.]}
}
{
}

\Verse{26}
{
    \zR{说 着，将要不理。}
}
{
    \eR{[English source not present in the checked-in Bencraft/Joly files.]}
}
{
}

\Verse{27}
{
    \zR{岫烟懊悔此来，知他脾气是这么着的，“一时我已说出，不好白回 去”。}
}
{
    \eR{[English source not present in the checked-in Bencraft/Joly files.]}
}
{
}

\Verse{28}
{
    \zR{又不好与他质证他会扶乩的话，只得陪着笑将袭人等性命关系的话说了一遍。}
}
{
    \eR{[English source not present in the checked-in Bencraft/Joly files.]}
}
{
}

\Verse{29}
{
    \zR{见妙玉略有活动，便起身拜了几拜。}
}
{
    \eR{[English source not present in the checked-in Bencraft/Joly files.]}
}
{
}

\Verse{30}
{
    \zR{妙玉叹道：“何必为人作嫁？}
}
{
    \eR{[English source not present in the checked-in Bencraft/Joly files.]}
}
{
}

\Verse{31}
{
    \zR{但是我进京以来， 素无人知，今日你来破例，恐将来缠绕不休。”}
}
{
    \eR{[English source not present in the checked-in Bencraft/Joly files.]}
}
{
}

\Verse{32}
{
    \zR{岫烟道：“我也一时不忍。}
}
{
    \eR{[English source not present in the checked-in Bencraft/Joly files.]}
}
{
}

\Verse{33}
{
    \zR{知你必 是慈悲的。}
}
{
    \eR{[English source not present in the checked-in Bencraft/Joly files.]}
}
{
}

\Verse{34}
{
    \zR{便是将来他人求你，愿不愿在你，谁敢相强？”}
}
{
    \eR{[English source not present in the checked-in Bencraft/Joly files.]}
}
{
}

\Verse{35}
{
    \zR{妙玉笑了一笑，叫道婆 焚香。}
}
{
    \eR{[English source not present in the checked-in Bencraft/Joly files.]}
}
{
}

\Verse{36}
{
    \zR{在箱子里找出沙盘乩架，书了符，命岫烟行礼祝告毕，起来同妙玉扶着乩。}
}
{
    \eR{[English source not present in the checked-in Bencraft/Joly files.]}
}
{
}

\Verse{37}
{
    \zR{不多时，只见那仙乩疾书道： 噫!来无迹，去无踪，青埂峰下倚古松。}
}
{
    \eR{[English source not present in the checked-in Bencraft/Joly files.]}
}
{
}

\Verse{38}
{
    \zR{欲追寻，山万重，入我门来一笑逢。}
}
{
    \eR{[English source not present in the checked-in Bencraft/Joly files.]}
}
{
}

\Verse{39}
{
    \zR{书毕，停了乩，岫烟便问：“请的是何仙？”}
}
{
    \eR{[English source not present in the checked-in Bencraft/Joly files.]}
}
{
}

\Verse{40}
{
    \zR{妙玉道：“请的是拐仙。”}
}
{
    \eR{[English source not present in the checked-in Bencraft/Joly files.]}
}
{
}

\Verse{41}
{
    \zR{岫烟录了 出来，请教妙玉识。}
}
{
    \eR{[English source not present in the checked-in Bencraft/Joly files.]}
}
{
}

\Verse{42}
{
    \zR{妙玉道：“这个可不能，连我也不懂。}
}
{
    \eR{[English source not present in the checked-in Bencraft/Joly files.]}
}
{
}

\Verse{43}
{
    \zR{你快拿去，他们的聪明 人多着哩。”}
}
{
    \eR{[English source not present in the checked-in Bencraft/Joly files.]}
}
{
}

\Verse{44}
{
    \zR{岫烟只得回来。}
}
{
    \eR{[English source not present in the checked-in Bencraft/Joly files.]}
}
{
}

\Verse{45}
{
    \zR{进入院中，各人都问：“怎么样了？”}
}
{
    \eR{[English source not present in the checked-in Bencraft/Joly files.]}
}
{
}

\Verse{46}
{
    \zR{岫烟不及细说，便将所录乩语递与李纨。}
}
{
    \eR{[English source not present in the checked-in Bencraft/Joly files.]}
}
{
}

\Verse{47}
{
    \zR{众姊妹及宝玉争看，都解的是：“一时要找是找不着的，然而丢是丢不了的，不知 几时不找便出来了。}
}
{
    \eR{[English source not present in the checked-in Bencraft/Joly files.]}
}
{
}

\Verse{48}
{
    \zR{但是青埂峰不知在那里？”}
}
{
    \eR{[English source not present in the checked-in Bencraft/Joly files.]}
}
{
}

\Verse{49}
{
    \zR{李纨道：“这是仙机隐语。}
}
{
    \eR{[English source not present in the checked-in Bencraft/Joly files.]}
}
{
}

\Verse{50}
{
    \zR{咱们家 里那里跑出青埂峰来？}
}
{
    \eR{[English source not present in the checked-in Bencraft/Joly files.]}
}
{
}

\Verse{51}
{
    \zR{必是谁怕查出，撂在有松树的山子石底下，也未可定。}
}
{
    \eR{[English source not present in the checked-in Bencraft/Joly files.]}
}
{
}

\Verse{52}
{
    \zR{独是 ‘入我门来’这句，到底是入谁的门呢？”}
}
{
    \eR{[English source not present in the checked-in Bencraft/Joly files.]}
}
{
}

\Verse{53}
{
    \zR{黛玉道：“不知请的是谁？”}
}
{
    \eR{[English source not present in the checked-in Bencraft/Joly files.]}
}
{
}

\Verse{54}
{
    \zR{岫烟道： “拐仙。”}
}
{
    \eR{[English source not present in the checked-in Bencraft/Joly files.]}
}
{
}

\Verse{55}
{
    \zR{探春道：“若是仙家的门，便难入了。”}
}
{
    \eR{[English source not present in the checked-in Bencraft/Joly files.]}
}
{
}

\Verse{56}
{
    \zR{袭人心里着忙，便捕风捉影的 混找，没一块石底下不找到，只是没有。}
}
{
    \eR{[English source not present in the checked-in Bencraft/Joly files.]}
}
{
}
\ChapterImage{img/chapters/191第95回 櫳翠庵妙玉扶乩玉.jpg}


\Verse{57}
{
    \zR{回到院中，宝玉也不问有无，只管傻笑。}
}
{
    \eR{[English source not present in the checked-in Bencraft/Joly files.]}
}
{
}

\Verse{58}
{
    \zR{麝月着急道：“小祖宗!你到底是那里丢的？}
}
{
    \eR{[English source not present in the checked-in Bencraft/Joly files.]}
}
{
}

\Verse{59}
{
    \zR{说明了，我们就是受罪，也在明处啊。”}
}
{
    \eR{[English source not present in the checked-in Bencraft/Joly files.]}
}
{
}

\Verse{60}
{
    \zR{宝玉笑道：“我说外头丢的，你们又不依。}
}
{
    \eR{[English source not present in the checked-in Bencraft/Joly files.]}
}
{
}

\Verse{61}
{
    \zR{你如今问我，我知道么？”}
}
{
    \eR{[English source not present in the checked-in Bencraft/Joly files.]}
}
{
}

\Verse{62}
{
    \zR{李纨探春道： “今儿从早起闹起，已到三更来的天了。}
}
{
    \eR{[English source not present in the checked-in Bencraft/Joly files.]}
}
{
}

\Verse{63}
{
    \zR{你瞧林妹妹已经掌不住，各自去了。}
}
{
    \eR{[English source not present in the checked-in Bencraft/Joly files.]}
}
{
}

\Verse{64}
{
    \zR{我们 也该歇歇儿了，明儿再闹罢。”}
}
{
    \eR{[English source not present in the checked-in Bencraft/Joly files.]}
}
{
}

\Verse{65}
{
    \zR{说着，大家散去。}
}
{
    \eR{[English source not present in the checked-in Bencraft/Joly files.]}
}
{
}

\Verse{66}
{
    \zR{宝玉即便睡下。}
}
{
    \eR{[English source not present in the checked-in Bencraft/Joly files.]}
}
{
}

\Verse{67}
{
    \zR{可怜袭人等哭一 回，想一回，一夜无眠，暂且不提。}
}
{
    \eR{[English source not present in the checked-in Bencraft/Joly files.]}
}
{
}

\Verse{68}
{
    \zR{且说黛玉先自回去，想起“金”“石”的旧话来，反自欢喜，心里也道：“和 尚道士的话真个信不得。}
}
{
    \eR{[English source not present in the checked-in Bencraft/Joly files.]}
}
{
}

\Verse{69}
{
    \zR{果真‘金’‘玉’有缘，宝玉如何能把这玉丢了呢？}
}
{
    \eR{[English source not present in the checked-in Bencraft/Joly files.]}
}
{
}

\Verse{70}
{
    \zR{或者 因我之事，拆散他们的‘金玉’，也未可知。”}
}
{
    \eR{[English source not present in the checked-in Bencraft/Joly files.]}
}
{
}

\Verse{71}
{
    \zR{想了半天，更觉安心，把这一天的 劳乏竟不理会，重新倒看起书来。}
}
{
    \eR{[English source not present in the checked-in Bencraft/Joly files.]}
}
{
}

\Verse{72}
{
    \zR{紫鹃倒觉身倦，连催黛玉睡下。}
}
{
    \eR{[English source not present in the checked-in Bencraft/Joly files.]}
}
{
}

\Verse{73}
{
    \zR{黛玉虽躺下，又 想到海棠花上，说：“这块玉原是胎里带来的，非比寻常之物，来去自有关系。}
}
{
    \eR{[English source not present in the checked-in Bencraft/Joly files.]}
}
{
}

\Verse{74}
{
    \zR{若 是这花主好事呢，不该失了这玉呀。}
}
{
    \eR{[English source not present in the checked-in Bencraft/Joly files.]}
}
{
}

\Verse{75}
{
    \zR{看来此花开的不祥，莫非他有不吉之事？”}
}
{
    \eR{[English source not present in the checked-in Bencraft/Joly files.]}
}
{
}

\Verse{76}
{
    \zR{不 觉又伤起心来。}
}
{
    \eR{[English source not present in the checked-in Bencraft/Joly files.]}
}
{
}

\Verse{77}
{
    \zR{又转想到喜事上头，此花又似应开，此玉又似应失：如此一悲一喜， 直想到五更方睡着。}
}
{
    \eR{[English source not present in the checked-in Bencraft/Joly files.]}
}
{
}

\Verse{78}
{
    \zR{次日，王夫人等早派人到当铺里去查问，凤姐暗中设法找寻。}
}
{
    \eR{[English source not present in the checked-in Bencraft/Joly files.]}
}
{
}

\Verse{79}
{
    \zR{一连闹了几天， 总无下落。}
}
{
    \eR{[English source not present in the checked-in Bencraft/Joly files.]}
}
{
}

\Verse{80}
{
    \zR{还喜贾母贾政未知。}
}
{
    \eR{[English source not present in the checked-in Bencraft/Joly files.]}
}
{
}

\Verse{81}
{
    \zR{袭人等每日提心吊胆。}
}
{
    \eR{[English source not present in the checked-in Bencraft/Joly files.]}
}
{
}

\Verse{82}
{
    \zR{宝玉也好几天不上学，只是 怔怔的，不言不语，没心没绪的。}
}
{
    \eR{[English source not present in the checked-in Bencraft/Joly files.]}
}
{
}

\Verse{83}
{
    \zR{王夫人只知他因失玉而起，也不大着意。}
}
{
    \eR{[English source not present in the checked-in Bencraft/Joly files.]}
}
{
}

\Verse{84}
{
    \zR{那日正 在纳闷，忽见贾琏进来请安，嘻嘻的笑道：“今日听得雨村打发人来告诉咱们二老
        爷，说舅太爷升了内阁大学士，奉旨来京，已定于明年正月二十日宣麻，有三百里 的文书去了。}
}
{
    \eR{[English source not present in the checked-in Bencraft/Joly files.]}
}
{
}

\Verse{85}
{
    \zR{想舅太爷昼夜趱行，半个多月就要到了。}
}
{
    \eR{[English source not present in the checked-in Bencraft/Joly files.]}
}
{
}

\Verse{86}
{
    \zR{侄儿特来回太太知道。”}
}
{
    \eR{[English source not present in the checked-in Bencraft/Joly files.]}
}
{
}

\Verse{87}
{
    \zR{王 夫人听说，便欢喜非常。}
}
{
    \eR{[English source not present in the checked-in Bencraft/Joly files.]}
}
{
}

\Verse{88}
{
    \zR{正想娘家人少，薛姨妈家又衰败了，兄弟又在外任照应不 着，今日忽听兄弟拜相回京，王家荣耀，将来宝玉都有倚靠，便把失玉的心又略放 开些了，天天专望兄弟来京。}
}
{
    \eR{[English source not present in the checked-in Bencraft/Joly files.]}
}
{
}

\Verse{89}
{
    \zR{忽一天，贾政进来，满脸泪痕，喘吁吁的说道：“你快去禀知老太太，即刻进 宫!不用多人的，是你伏侍进去。}
}
{
    \eR{[English source not present in the checked-in Bencraft/Joly files.]}
}
{
}

\Verse{90}
{
    \zR{因娘娘忽得暴病，现在太监在外立等。}
}
{
    \eR{[English source not present in the checked-in Bencraft/Joly files.]}
}
{
}

\Verse{91}
{
    \zR{他说：‘太 医院已经奏明痰厥，不能医治。’”}
}
{
    \eR{[English source not present in the checked-in Bencraft/Joly files.]}
}
{
}

\Verse{92}
{
    \zR{王夫人听说，便大哭起来。}
}
{
    \eR{[English source not present in the checked-in Bencraft/Joly files.]}
}
{
}

\Verse{93}
{
    \zR{贾政道：“这不是 哭的时候，快快去请老太太。}
}
{
    \eR{[English source not present in the checked-in Bencraft/Joly files.]}
}
{
}

\Verse{94}
{
    \zR{说得宽缓些，不要吓坏了老人家。”}
}
{
    \eR{[English source not present in the checked-in Bencraft/Joly files.]}
}
{
}

\Verse{95}
{
    \zR{贾政说着，出来 吩咐家人伺候。}
}
{
    \eR{[English source not present in the checked-in Bencraft/Joly files.]}
}
{
}

\Verse{96}
{
    \zR{王夫人收了泪，去请贾母，只说元妃有病，进去请安。}
}
{
    \eR{[English source not present in the checked-in Bencraft/Joly files.]}
}
{
}

\Verse{97}
{
    \zR{贾母念佛道： “怎么又病了？}
}
{
    \eR{[English source not present in the checked-in Bencraft/Joly files.]}
}
{
}

\Verse{98}
{
    \zR{前番吓的我了不得，后来又打听错了。}
}
{
    \eR{[English source not present in the checked-in Bencraft/Joly files.]}
}
{
}

\Verse{99}
{
    \zR{这回情愿再错了也罢。”}
}
{
    \eR{[English source not present in the checked-in Bencraft/Joly files.]}
}
{
}

\Verse{100}
{
    \zR{王 夫人一面回答，一面催鸳鸯等开箱取衣饰穿戴起来。}
}
{
    \eR{[English source not present in the checked-in Bencraft/Joly files.]}
}
{
}

\Verse{101}
{
    \zR{王夫人赶着回到自己房中，也 穿戴好了，过来伺候。}
}
{
    \eR{[English source not present in the checked-in Bencraft/Joly files.]}
}
{
}

\Verse{102}
{
    \zR{一时出厅，上轿进宫不提。}
}
{
    \eR{[English source not present in the checked-in Bencraft/Joly files.]}
}
{
}

\Verse{103}
{
    \zR{且说元春自选了凤藻宫后，圣眷隆重，身体发福，未免举动费力。}
}
{
    \eR{[English source not present in the checked-in Bencraft/Joly files.]}
}
{
}

\Verse{104}
{
    \zR{每日起居劳 乏，时发痰疾。}
}
{
    \eR{[English source not present in the checked-in Bencraft/Joly files.]}
}
{
}

\Verse{105}
{
    \zR{因前日侍宴回宫，偶沾寒气，勾起旧病。}
}
{
    \eR{[English source not present in the checked-in Bencraft/Joly files.]}
}
{
}

\Verse{106}
{
    \zR{不料此回甚属利害，竟至 痰气壅塞，四肢厥冷。}
}
{
    \eR{[English source not present in the checked-in Bencraft/Joly files.]}
}
{
}

\Verse{107}
{
    \zR{一面奏明，即召太医调治。}
}
{
    \eR{[English source not present in the checked-in Bencraft/Joly files.]}
}
{
}

\Verse{108}
{
    \zR{岂知汤药不进，连用通关之剂， 并不见效。}
}
{
    \eR{[English source not present in the checked-in Bencraft/Joly files.]}
}
{
}

\Verse{109}
{
    \zR{内官忧虑，奏请预办后事，所以传旨命贾氏椒房进见。}
}
{
    \eR{[English source not present in the checked-in Bencraft/Joly files.]}
}
{
}

\Verse{110}
{
    \zR{贾母王夫人遵旨 进宫，见元妃痰塞口涎，不能言语。}
}
{
    \eR{[English source not present in the checked-in Bencraft/Joly files.]}
}
{
}

\Verse{111}
{
    \zR{见了贾母，只有悲泣之状，却没眼泪。}
}
{
    \eR{[English source not present in the checked-in Bencraft/Joly files.]}
}
{
}
\ChapterImage{img/chapters/192第95回 因訛成寶元妃薨逝.jpg}


\Verse{112}
{
    \zR{贾母进 前请安，奏些宽慰的话。}
}
{
    \eR{[English source not present in the checked-in Bencraft/Joly files.]}
}
{
}

\Verse{113}
{
    \zR{少时贾政等职名递进，宫嫔传奏，元妃目不能顾，渐渐脸 色改变。}
}
{
    \eR{[English source not present in the checked-in Bencraft/Joly files.]}
}
{
}

\Verse{114}
{
    \zR{内官太监即要奏闻，恐派各妃看视，椒房姻戚未便久羁，请在外宫伺候。}
}
{
    \eR{[English source not present in the checked-in Bencraft/Joly files.]}
}
{
}

\Verse{115}
{
    \zR{贾母王夫人怎忍便离，无奈国家制度，只得下来，又不敢啼哭，惟有心内悲感。}
}
{
    \eR{[English source not present in the checked-in Bencraft/Joly files.]}
}
{
}

\Verse{116}
{
    \zR{朝门内官员有信。}
}
{
    \eR{[English source not present in the checked-in Bencraft/Joly files.]}
}
{
}

\Verse{117}
{
    \zR{不多时，只见太监出来，立传钦天监。}
}
{
    \eR{[English source not present in the checked-in Bencraft/Joly files.]}
}
{
}

\Verse{118}
{
    \zR{贾母便知不好，尚未 敢动。}
}
{
    \eR{[English source not present in the checked-in Bencraft/Joly files.]}
}
{
}

\Verse{119}
{
    \zR{稍刻，小太监传谕出来，说：“贾娘娘薨逝。”}
}
{
    \eR{[English source not present in the checked-in Bencraft/Joly files.]}
}
{
}

\Verse{120}
{
    \zR{是年甲寅年十二月十八日立 春，元妃薨日，是十二月十九日，已交卯年寅月，存年四十三岁。}
}
{
    \eR{[English source not present in the checked-in Bencraft/Joly files.]}
}
{
}

\Verse{121}
{
    \zR{贾母含悲起身， 只得出宫上轿回家。}
}
{
    \eR{[English source not present in the checked-in Bencraft/Joly files.]}
}
{
}

\Verse{122}
{
    \zR{贾政等亦已得信，一路悲戚。}
}
{
    \eR{[English source not present in the checked-in Bencraft/Joly files.]}
}
{
}

\Verse{123}
{
    \zR{到家中，邢夫人、李纨、凤姐、 宝玉等出厅，分东西迎着贾母，请了安，并贾政王夫人请安，大家哭泣不提。}
}
{
    \eR{[English source not present in the checked-in Bencraft/Joly files.]}
}
{
}

\Verse{124}
{
    \zR{次日早起，凡有品级的，按贵妃丧礼进内请安哭临。}
}
{
    \eR{[English source not present in the checked-in Bencraft/Joly files.]}
}
{
}

\Verse{125}
{
    \zR{贾政又是工部，虽按照仪 注办理，未免堂上又要周旋他些，同事又要请教他，所以两头更忙，非比从前太后 与周妃的丧事了。}
}
{
    \eR{[English source not present in the checked-in Bencraft/Joly files.]}
}
{
}

\Verse{126}
{
    \zR{但元妃并无所出，惟谥曰贤淑贵妃。}
}
{
    \eR{[English source not present in the checked-in Bencraft/Joly files.]}
}
{
}

\Verse{127}
{
    \zR{此是王家制度，不必多赘。}
}
{
    \eR{[English source not present in the checked-in Bencraft/Joly files.]}
}
{
}

\Verse{128}
{
    \zR{只讲贾府中男女，天天进宫，忙的了不得。}
}
{
    \eR{[English source not present in the checked-in Bencraft/Joly files.]}
}
{
}

\Verse{129}
{
    \zR{幸喜凤姐儿近日身子好些，还得出来照 应家事，又要预备王子腾进京，接风贺喜。}
}
{
    \eR{[English source not present in the checked-in Bencraft/Joly files.]}
}
{
}

\Verse{130}
{
    \zR{凤姐胞兄王仁，知道叔叔入了内阁，仍 带家眷来京。}
}
{
    \eR{[English source not present in the checked-in Bencraft/Joly files.]}
}
{
}

\Verse{131}
{
    \zR{凤姐心里喜欢，便有些心病，有这些娘家的人也便撂开，所以身子倒 觉比先好了些。}
}
{
    \eR{[English source not present in the checked-in Bencraft/Joly files.]}
}
{
}

\Verse{132}
{
    \zR{王夫人看见凤姐照旧办事，又把担子卸了一半，又眼见兄弟来京， 诸事放心，倒觉安静些。}
}
{
    \eR{[English source not present in the checked-in Bencraft/Joly files.]}
}
{
}

\Verse{133}
{
    \zR{独有宝玉原是无职之人，又不念书，代儒学里知他家里有事，也不来管他；贾 政正忙，自然没有空儿查他。}
}
{
    \eR{[English source not present in the checked-in Bencraft/Joly files.]}
}
{
}

\Verse{134}
{
    \zR{想来宝玉趁此机会，竟可与姊妹们天天畅乐；不料他 自失了玉后，终日懒怠走动，说话也糊涂了。}
}
{
    \eR{[English source not present in the checked-in Bencraft/Joly files.]}
}
{
}

\Verse{135}
{
    \zR{并贾母等出门回来，有人叫他去请安， 便去；没人叫他，他也不动。}
}
{
    \eR{[English source not present in the checked-in Bencraft/Joly files.]}
}
{
}

\Verse{136}
{
    \zR{袭人等怀着鬼胎，又不敢去招惹他，恐他生气。}
}
{
    \eR{[English source not present in the checked-in Bencraft/Joly files.]}
}
{
}

\Verse{137}
{
    \zR{每天 茶饭，端到面前便吃，不来也不要。}
}
{
    \eR{[English source not present in the checked-in Bencraft/Joly files.]}
}
{
}

\Verse{138}
{
    \zR{袭人看这光景，不像是有气，竟像是有病的。}
}
{
    \eR{[English source not present in the checked-in Bencraft/Joly files.]}
}
{
}

\Verse{139}
{
    \zR{袭人偷着空儿到潇湘馆告诉紫鹃，说是：“二爷这么着，求姑娘给他开导开导。”}
}
{
    \eR{[English source not present in the checked-in Bencraft/Joly files.]}
}
{
}

\Verse{140}
{
    \zR{紫鹃虽即告诉黛玉，只因黛玉想着亲事上头，一定是自己了，如今见了他，反觉不 好意思：“若是他来呢，原是小时在一处的，也难不理他；若说我去找他，断断使
        不得。”}
}
{
    \eR{[English source not present in the checked-in Bencraft/Joly files.]}
}
{
}

\Verse{141}
{
    \zR{所以黛玉不肯过来。}
}
{
    \eR{[English source not present in the checked-in Bencraft/Joly files.]}
}
{
}

\Verse{142}
{
    \zR{袭人又背地里去告诉探春。}
}
{
    \eR{[English source not present in the checked-in Bencraft/Joly files.]}
}
{
}

\Verse{143}
{
    \zR{那知探春心里明明知道海 棠开得怪异，“宝玉”失的更奇，接连着元妃姐姐薨逝，谅家道不祥，日日愁闷， 那有心肠去劝宝玉？}
}
{
    \eR{[English source not present in the checked-in Bencraft/Joly files.]}
}
{
}

\Verse{144}
{
    \zR{况兄妹们男女有别，只好过来一两次，宝玉又终是懒懒的，所 以也不大常来。}
}
{
    \eR{[English source not present in the checked-in Bencraft/Joly files.]}
}
{
}

\Verse{145}
{
    \zR{宝钗也知失玉。}
}
{
    \eR{[English source not present in the checked-in Bencraft/Joly files.]}
}
{
}

\Verse{146}
{
    \zR{因薛姨妈那日应了宝玉的亲事，回去便告诉了宝钗。}
}
{
    \eR{[English source not present in the checked-in Bencraft/Joly files.]}
}
{
}

\Verse{147}
{
    \zR{薛姨妈还 说：“虽是你姨妈说了，我还没有应准，说等你哥哥回来再定。}
}
{
    \eR{[English source not present in the checked-in Bencraft/Joly files.]}
}
{
}

\Verse{148}
{
    \zR{你愿意不愿意？”}
}
{
    \eR{[English source not present in the checked-in Bencraft/Joly files.]}
}
{
}

\Verse{149}
{
    \zR{宝钗反正色的对母亲道：“妈妈这话说错了。}
}
{
    \eR{[English source not present in the checked-in Bencraft/Joly files.]}
}
{
}

\Verse{150}
{
    \zR{女孩儿家的事情是父母作主的，如今 我父亲没了，妈妈应该作主的，再不然问哥哥。}
}
{
    \eR{[English source not present in the checked-in Bencraft/Joly files.]}
}
{
}

\Verse{151}
{
    \zR{怎么问起我来？”}
}
{
    \eR{[English source not present in the checked-in Bencraft/Joly files.]}
}
{
}

\Verse{152}
{
    \zR{所以薛姨妈更爱 惜他，说他虽是从小娇养惯的，却也生来的贞静，因此在他面前反不提起宝玉了。}
}
{
    \eR{[English source not present in the checked-in Bencraft/Joly files.]}
}
{
}

\Verse{153}
{
    \zR{宝钗自从听此一说，把“宝玉”两字自然更不提起了。}
}
{
    \eR{[English source not present in the checked-in Bencraft/Joly files.]}
}
{
}

\Verse{154}
{
    \zR{如今虽然听见失了玉，心里 也甚惊疑，倒不好问，只得听旁人说去，竟像不与自己相干的。}
}
{
    \eR{[English source not present in the checked-in Bencraft/Joly files.]}
}
{
}

\Verse{155}
{
    \zR{只有薛姨妈打发丫 头过来了好几次问信。}
}
{
    \eR{[English source not present in the checked-in Bencraft/Joly files.]}
}
{
}

\Verse{156}
{
    \zR{因他自己的儿子薛蟠的事焦心，只等哥哥进京，便好为他出 脱罪名；又知元妃已薨，虽然贾府忙乱，却得凤姐好了，出来理家，所以也不大过 这边来。}
}
{
    \eR{[English source not present in the checked-in Bencraft/Joly files.]}
}
{
}

\Verse{157}
{
    \zR{这里只苦了袭人，在宝玉跟前低声下气的伏侍劝慰，宝玉竟是不懂。}
}
{
    \eR{[English source not present in the checked-in Bencraft/Joly files.]}
}
{
}

\Verse{158}
{
    \zR{袭人 只有暗暗的着急而已。}
}
{
    \eR{[English source not present in the checked-in Bencraft/Joly files.]}
}
{
}

\Verse{159}
{
    \zR{过了几日，元妃停灵寝庙，贾母等送殡去了几天。}
}
{
    \eR{[English source not present in the checked-in Bencraft/Joly files.]}
}
{
}

\Verse{160}
{
    \zR{岂知宝玉一日呆似一日，也 不发烧，也不疼痛，只是吃不像吃，睡不像睡，甚至说话都无头绪。}
}
{
    \eR{[English source not present in the checked-in Bencraft/Joly files.]}
}
{
}

\Verse{161}
{
    \zR{那袭人麝月等 一发慌了，回过凤姐几次。}
}
{
    \eR{[English source not present in the checked-in Bencraft/Joly files.]}
}
{
}

\Verse{162}
{
    \zR{凤姐不时过来。}
}
{
    \eR{[English source not present in the checked-in Bencraft/Joly files.]}
}
{
}

\Verse{163}
{
    \zR{起先道是找不着玉生气，如今看他失魂 落魄的样子，只有日日请医调治。}
}
{
    \eR{[English source not present in the checked-in Bencraft/Joly files.]}
}
{
}

\Verse{164}
{
    \zR{煎药吃了好几剂，只有添病的，没有减病的。}
}
{
    \eR{[English source not present in the checked-in Bencraft/Joly files.]}
}
{
}

\Verse{165}
{
    \zR{及 至问他那里不舒服，宝玉也不说出来。}
}
{
    \eR{[English source not present in the checked-in Bencraft/Joly files.]}
}
{
}

\Verse{166}
{
    \zR{直至元妃事毕，贾母惦记宝玉，亲自到园看 视，王夫人也随过来。}
}
{
    \eR{[English source not present in the checked-in Bencraft/Joly files.]}
}
{
}

\Verse{167}
{
    \zR{袭人等叫宝玉接出去请安。}
}
{
    \eR{[English source not present in the checked-in Bencraft/Joly files.]}
}
{
}
\ChapterImage{img/chapters/193第95-96回 以假混真寶玉癡癲.jpg}


\Verse{168}
{
    \zR{宝玉虽说是病，每日原起来行动， 今日叫他接贾母去，他依然仍是请安，惟是袭人在旁扶着指教。}
}
{
    \eR{[English source not present in the checked-in Bencraft/Joly files.]}
}
{
}

\Verse{169}
{
    \zR{贾母见了，便道： “我的儿，我打量你怎么病着，故此过来瞧你。}
}
{
    \eR{[English source not present in the checked-in Bencraft/Joly files.]}
}
{
}

\Verse{170}
{
    \zR{今你依旧的模样儿，我的心放了好 些。”}
}
{
    \eR{[English source not present in the checked-in Bencraft/Joly files.]}
}
{
}

\Verse{171}
{
    \zR{王夫人也自然是宽心的。}
}
{
    \eR{[English source not present in the checked-in Bencraft/Joly files.]}
}
{
}

\Verse{172}
{
    \zR{但宝玉并不回答，只管嘻嘻的笑。}
}
{
    \eR{[English source not present in the checked-in Bencraft/Joly files.]}
}
{
}

\Verse{173}
{
    \zR{贾母等进屋坐下， 问他的话，袭人教一句，他说一句，大不似往常，直是一个傻子似的。}
}
{
    \eR{[English source not present in the checked-in Bencraft/Joly files.]}
}
{
}

\Verse{174}
{
    \zR{贾母愈看愈 疑，便说：“我才进来看时，不见有什么病；如今细细一瞧，这病果然不轻，竟是 神魂失散的样子。}
}
{
    \eR{[English source not present in the checked-in Bencraft/Joly files.]}
}
{
}

\Verse{175}
{
    \zR{到底因什么起的呢？”}
}
{
    \eR{[English source not present in the checked-in Bencraft/Joly files.]}
}
{
}

\Verse{176}
{
    \zR{王夫人知事难瞒，又瞧瞧袭人怪可怜的样 子，只得便依着宝玉先前的话，将那往临安伯府里去听戏时丢了这块玉的话悄悄的 告诉了一遍，心里也？}
}
{
    \eR{[English source not present in the checked-in Bencraft/Joly files.]}
}
{
}

\Verse{177}
{
    \zR{徨的很，生恐贾母着急。}
}
{
    \eR{[English source not present in the checked-in Bencraft/Joly files.]}
}
{
}

\Verse{178}
{
    \zR{并说：“现在着人在四下里找寻。}
}
{
    \eR{[English source not present in the checked-in Bencraft/Joly files.]}
}
{
}

\Verse{179}
{
    \zR{求签问卦，都说在当铺里找，少不得找着的。”}
}
{
    \eR{[English source not present in the checked-in Bencraft/Joly files.]}
}
{
}

\Verse{180}
{
    \zR{贾母听了，急得站起来，眼泪直流， 说道：“这件玉如何是丢得的!你们忒不懂事了!难道老爷也是撂开手的不成？”}
}
{
    \eR{[English source not present in the checked-in Bencraft/Joly files.]}
}
{
}

\Verse{181}
{
    \zR{王 夫人知贾母生气，叫袭人等跪下，自己敛容低首回说：“媳妇恐老太太着急，老爷 生气，都没敢回。”}
}
{
    \eR{[English source not present in the checked-in Bencraft/Joly files.]}
}
{
}

\Verse{182}
{
    \zR{贾母咳道：“这是宝玉的命根子，因丢了，所以他这么失魂丧 魄的。}
}
{
    \eR{[English source not present in the checked-in Bencraft/Joly files.]}
}
{
}

\Verse{183}
{
    \zR{还了得!这玉是满城里都知道的，谁检了去，肯叫你们找出来么？}
}
{
    \eR{[English source not present in the checked-in Bencraft/Joly files.]}
}
{
}

\Verse{184}
{
    \zR{叫人快快请 老爷，我与他说。”}
}
{
    \eR{[English source not present in the checked-in Bencraft/Joly files.]}
}
{
}

\Verse{185}
{
    \zR{那时吓得王夫人袭人等俱哀告道：“老太太这一生气，回来老 爷更了不得了。}
}
{
    \eR{[English source not present in the checked-in Bencraft/Joly files.]}
}
{
}

\Verse{186}
{
    \zR{现在宝玉病着，交给我们尽命的找来就是了。”}
}
{
    \eR{[English source not present in the checked-in Bencraft/Joly files.]}
}
{
}

\Verse{187}
{
    \zR{贾母道：“你们怕 老爷生气，有我呢。”}
}
{
    \eR{[English source not present in the checked-in Bencraft/Joly files.]}
}
{
}

\Verse{188}
{
    \zR{便叫麝月传人去请。}
}
{
    \eR{[English source not present in the checked-in Bencraft/Joly files.]}
}
{
}

\Verse{189}
{
    \zR{不一时传话进来，说：“老爷谢客去了。”}
}
{
    \eR{[English source not present in the checked-in Bencraft/Joly files.]}
}
{
}

\Verse{190}
{
    \zR{贾母道：“不用他也使得。}
}
{
    \eR{[English source not present in the checked-in Bencraft/Joly files.]}
}
{
}

\Verse{191}
{
    \zR{你们便 说我说的话，暂且也不用责罚下人。}
}
{
    \eR{[English source not present in the checked-in Bencraft/Joly files.]}
}
{
}

\Verse{192}
{
    \zR{我便叫琏儿来，写出赏格，悬在前日经过的地 方，便说：‘有人检得送来者，情愿送银一万两；如有知人检得，送信找得者，送 银五千两。’}
}
{
    \eR{[English source not present in the checked-in Bencraft/Joly files.]}
}
{
}

\Verse{193}
{
    \zR{如真有了，不可吝惜银子。}
}
{
    \eR{[English source not present in the checked-in Bencraft/Joly files.]}
}
{
}

\Verse{194}
{
    \zR{这么一找，少不得就找出来了。}
}
{
    \eR{[English source not present in the checked-in Bencraft/Joly files.]}
}
{
}

\Verse{195}
{
    \zR{若是靠着 咱们家几个人找，就找一辈子也不能得！”}
}
{
    \eR{[English source not present in the checked-in Bencraft/Joly files.]}
}
{
}

\Verse{196}
{
    \zR{王夫人也不敢直言。}
}
{
    \eR{[English source not present in the checked-in Bencraft/Joly files.]}
}
{
}

\Verse{197}
{
    \zR{贾母传话告诉贾琏， 叫他速办去了。}
}
{
    \eR{[English source not present in the checked-in Bencraft/Joly files.]}
}
{
}

\Verse{198}
{
    \zR{贾母便叫人：“将宝玉动用之物，都搬到我那里去。}
}
{
    \eR{[English source not present in the checked-in Bencraft/Joly files.]}
}
{
}

\Verse{199}
{
    \zR{只派袭人秋纹 跟过来，馀者仍留园内看屋子。”}
}
{
    \eR{[English source not present in the checked-in Bencraft/Joly files.]}
}
{
}

\Verse{200}
{
    \zR{宝玉听了，总不言语，只是傻笑。}
}
{
    \eR{[English source not present in the checked-in Bencraft/Joly files.]}
}
{
}

\Verse{201}
{
    \zR{贾母便携了宝 玉起身，袭人等搀扶出园。}
}
{
    \eR{[English source not present in the checked-in Bencraft/Joly files.]}
}
{
}

\Verse{202}
{
    \zR{回到自己房中，叫王夫人坐下，看人收拾里间屋内安置，便对王夫人道：“你 知道我的意思么？}
}
{
    \eR{[English source not present in the checked-in Bencraft/Joly files.]}
}
{
}

\Verse{203}
{
    \zR{我为的是园里人少，怡红院的花树忽萎忽开，有些奇怪。}
}
{
    \eR{[English source not present in the checked-in Bencraft/Joly files.]}
}
{
}

\Verse{204}
{
    \zR{头里仗 着那块玉能除邪祟，如今玉丢了，只怕邪气易侵，所以我带过他来一块儿住着。}
}
{
    \eR{[English source not present in the checked-in Bencraft/Joly files.]}
}
{
}

\Verse{205}
{
    \zR{这 几天也不用叫他出去。}
}
{
    \eR{[English source not present in the checked-in Bencraft/Joly files.]}
}
{
}

\Verse{206}
{
    \zR{大夫来，就在这里瞧。”}
}
{
    \eR{[English source not present in the checked-in Bencraft/Joly files.]}
}
{
}

\Verse{207}
{
    \zR{王夫人听说，便接口道：“老太太 想的自然是。}
}
{
    \eR{[English source not present in the checked-in Bencraft/Joly files.]}
}
{
}

\Verse{208}
{
    \zR{如今宝玉同着老太太住了，老太太的福气大，不论什么都压住了。”}
}
{
    \eR{[English source not present in the checked-in Bencraft/Joly files.]}
}
{
}

\Verse{209}
{
    \zR{贾母道：“什么福气!不过我屋里干净些，经卷也多，都可以念念，定定心神。}
}
{
    \eR{[English source not present in the checked-in Bencraft/Joly files.]}
}
{
}

\Verse{210}
{
    \zR{你 问宝玉好不好？”}
}
{
    \eR{[English source not present in the checked-in Bencraft/Joly files.]}
}
{
}

\Verse{211}
{
    \zR{那宝玉见问只是笑。}
}
{
    \eR{[English source not present in the checked-in Bencraft/Joly files.]}
}
{
}

\Verse{212}
{
    \zR{袭人叫他说好，宝玉也就说好。}
}
{
    \eR{[English source not present in the checked-in Bencraft/Joly files.]}
}
{
}

\Verse{213}
{
    \zR{王夫人见了 这般光景，未免落泪，在贾母这里，不敢出声。}
}
{
    \eR{[English source not present in the checked-in Bencraft/Joly files.]}
}
{
}

\Verse{214}
{
    \zR{贾母知王夫人着急，便说道：“你 回去罢，这里有我调停他。}
}
{
    \eR{[English source not present in the checked-in Bencraft/Joly files.]}
}
{
}

\Verse{215}
{
    \zR{晚上老爷回来，告诉他不必来见我，不许言语就是了。”}
}
{
    \eR{[English source not present in the checked-in Bencraft/Joly files.]}
}
{
}

\Verse{216}
{
    \zR{王夫人去后，贾母叫鸳鸯找些安神定魄的药，按方吃了，不提。}
}
{
    \eR{[English source not present in the checked-in Bencraft/Joly files.]}
}
{
}

\Verse{217}
{
    \zR{且说贾政当晚回家，在车内听见道儿上人说道：“人要发财，也容易的很。”}
}
{
    \eR{[English source not present in the checked-in Bencraft/Joly files.]}
}
{
}

\Verse{218}
{
    \zR{那个问道：“怎么见得？”}
}
{
    \eR{[English source not present in the checked-in Bencraft/Joly files.]}
}
{
}

\Verse{219}
{
    \zR{这个人又道：“今日听见荣府里丢了什么哥儿的玉了， 贴着招帖儿，上头写着玉的大小式样颜色，说有人检了送去，就给一万两银子。}
}
{
    \eR{[English source not present in the checked-in Bencraft/Joly files.]}
}
{
}

\Verse{220}
{
    \zR{送 信的还给五千呢。”}
}
{
    \eR{[English source not present in the checked-in Bencraft/Joly files.]}
}
{
}

\Verse{221}
{
    \zR{贾政虽未听得如此真切，心里诧异，急忙赶回，便叫门上的人， 问起那事来。}
}
{
    \eR{[English source not present in the checked-in Bencraft/Joly files.]}
}
{
}

\Verse{222}
{
    \zR{门上的人禀道：“奴才头里也不知道，今儿晌午琏二爷传出老太太的 话，叫人去贴帖儿，才知道的。”}
}
{
    \eR{[English source not present in the checked-in Bencraft/Joly files.]}
}
{
}
\ChapterImage{img/chapters/194第95-96回 瞞消息鳳姐設奇謀.jpg}


\Verse{223}
{
    \zR{贾政便叹气道：“家道该衰!偏生养这么一个孽 障!才养他的时候，满街的谣言，隔了十几年略好了些。}
}
{
    \eR{[English source not present in the checked-in Bencraft/Joly files.]}
}
{
}

\Verse{224}
{
    \zR{这会子又大张晓谕的找玉， 成何道理！”}
}
{
    \eR{[English source not present in the checked-in Bencraft/Joly files.]}
}
{
}

\Verse{225}
{
    \zR{说着，忙走进里头去问王夫人。}
}
{
    \eR{[English source not present in the checked-in Bencraft/Joly files.]}
}
{
}

\Verse{226}
{
    \zR{王夫人便一五一十的告诉。}
}
{
    \eR{[English source not present in the checked-in Bencraft/Joly files.]}
}
{
}

\Verse{227}
{
    \zR{贾政知是 老太太的主意，又不敢违拗，只抱怨王夫人几句。}
}
{
    \eR{[English source not present in the checked-in Bencraft/Joly files.]}
}
{
}

\Verse{228}
{
    \zR{又走出来，叫瞒着老太太，背地 里揭了这个帖儿下来。}
}
{
    \eR{[English source not present in the checked-in Bencraft/Joly files.]}
}
{
}

\Verse{229}
{
    \zR{岂知早有那些游手好闲的人揭了去了。}
}
{
    \eR{[English source not present in the checked-in Bencraft/Joly files.]}
}
{
}

\Verse{230}
{
    \zR{过了些时，竟有人到荣府门上，口称送玉来的。}
}
{
    \eR{[English source not present in the checked-in Bencraft/Joly files.]}
}
{
}

\Verse{231}
{
    \zR{家人们听见，喜欢的了不得， 便说：“拿来，我给你回去。”}
}
{
    \eR{[English source not present in the checked-in Bencraft/Joly files.]}
}
{
}

\Verse{232}
{
    \zR{那人便怀内掏出赏格来，指给门上的人瞧，说：“这 不是你们府上的帖子？}
}
{
    \eR{[English source not present in the checked-in Bencraft/Joly files.]}
}
{
}

\Verse{233}
{
    \zR{写明送玉的给银一万两。}
}
{
    \eR{[English source not present in the checked-in Bencraft/Joly files.]}
}
{
}

\Verse{234}
{
    \zR{二太爷，你们这会子瞧我穷，回来 我得了银子，就是财主了，别这么待理不理的。”}
}
{
    \eR{[English source not present in the checked-in Bencraft/Joly files.]}
}
{
}

\Verse{235}
{
    \zR{门上人听他的话头儿硬，便说道： “你到底略给我瞧瞧，我好给你回。”}
}
{
    \eR{[English source not present in the checked-in Bencraft/Joly files.]}
}
{
}

\Verse{236}
{
    \zR{那人初倒不肯，后来听人说得有理，便掏出 那玉，托在掌中一扬，说：“这是不是？”}
}
{
    \eR{[English source not present in the checked-in Bencraft/Joly files.]}
}
{
}

\Verse{237}
{
    \zR{众家人原是在外服役，只知有玉，也不 常见，今日才看见这玉的模样儿了，急忙跑到里头抢头报的似的。}
}
{
    \eR{[English source not present in the checked-in Bencraft/Joly files.]}
}
{
}

\Verse{238}
{
    \zR{那日贾政贾赦出 门，只有贾琏在家。}
}
{
    \eR{[English source not present in the checked-in Bencraft/Joly files.]}
}
{
}

\Verse{239}
{
    \zR{众人回明，贾琏还问：“真不真？”}
}
{
    \eR{[English source not present in the checked-in Bencraft/Joly files.]}
}
{
}

\Verse{240}
{
    \zR{门上人口称：“亲眼见过， 只是不给奴才，要见主子，一手交银，一手交玉。”}
}
{
    \eR{[English source not present in the checked-in Bencraft/Joly files.]}
}
{
}

\Verse{241}
{
    \zR{贾琏却也喜欢，忙去禀知王夫 人，即便回明贾母，把个袭人乐的合掌念佛。}
}
{
    \eR{[English source not present in the checked-in Bencraft/Joly files.]}
}
{
}

\Verse{242}
{
    \zR{贾母并不改口，一叠连声：“快叫琏 儿请那人到书房里坐着，将玉取来一看，即便给银。”}
}
{
    \eR{[English source not present in the checked-in Bencraft/Joly files.]}
}
{
}

\Verse{243}
{
    \zR{贾琏依言，请那人进来，当 客待他，用好言道谢：“要借这玉送到里头本人见了，谢银分厘不短。”}
}
{
    \eR{[English source not present in the checked-in Bencraft/Joly files.]}
}
{
}

\Verse{244}
{
    \zR{那人只得 将一个红绸子包儿送过去。}
}
{
    \eR{[English source not present in the checked-in Bencraft/Joly files.]}
}
{
}

\Verse{245}
{
    \zR{贾琏打开一看，可不是那一块晶莹美玉吗？}
}
{
    \eR{[English source not present in the checked-in Bencraft/Joly files.]}
}
{
}

\Verse{246}
{
    \zR{贾琏素昔原 不理论，今日倒要看看。}
}
{
    \eR{[English source not present in the checked-in Bencraft/Joly files.]}
}
{
}

\Verse{247}
{
    \zR{看了半日，上面的字也仿佛认得出来，什么“除邪祟”等 字。}
}
{
    \eR{[English source not present in the checked-in Bencraft/Joly files.]}
}
{
}

\Verse{248}
{
    \zR{贾琏看了，喜之不胜，便叫家人伺候，忙忙的送与贾母王夫人认去。}
}
{
    \eR{[English source not present in the checked-in Bencraft/Joly files.]}
}
{
}

\Verse{249}
{
    \zR{这会子惊动了合家的人，都等着争看。}
}
{
    \eR{[English source not present in the checked-in Bencraft/Joly files.]}
}
{
}

\Verse{250}
{
    \zR{凤姐见贾琏进来，便劈手夺去，不敢先 看，送到贾母手里，贾琏笑道：“你这么一点儿事，还不叫我献功呢。”}
}
{
    \eR{[English source not present in the checked-in Bencraft/Joly files.]}
}
{
}

\Verse{251}
{
    \zR{贾母打开 看时，只见那玉比先前昏暗了好些，一面用手擦摸，鸳鸯拿上眼镜儿来，戴着一瞧， 说：“奇怪。}
}
{
    \eR{[English source not present in the checked-in Bencraft/Joly files.]}
}
{
}

\Verse{252}
{
    \zR{这块玉倒是的，怎么把头里的宝色都没了呢？”}
}
{
    \eR{[English source not present in the checked-in Bencraft/Joly files.]}
}
{
}

\Verse{253}
{
    \zR{王夫人看了一会子， 也认不出，便叫凤姐过来看。}
}
{
    \eR{[English source not present in the checked-in Bencraft/Joly files.]}
}
{
}

\Verse{254}
{
    \zR{凤姐看了道：“像倒像，只是颜色不大对，不如叫宝 兄弟自己一看，就知道了。”}
}
{
    \eR{[English source not present in the checked-in Bencraft/Joly files.]}
}
{
}

\Verse{255}
{
    \zR{袭人在旁，也看着未必是那一块，只是盼得的心盛， 也不敢说出不像来。}
}
{
    \eR{[English source not present in the checked-in Bencraft/Joly files.]}
}
{
}

\Verse{256}
{
    \zR{凤姐于是从贾母手中接过来，同着袭人，拿来给宝玉瞧。}
}
{
    \eR{[English source not present in the checked-in Bencraft/Joly files.]}
}
{
}

\Verse{257}
{
    \zR{这时 宝玉正睡着才醒。}
}
{
    \eR{[English source not present in the checked-in Bencraft/Joly files.]}
}
{
}

\Verse{258}
{
    \zR{凤姐告诉道：“你的玉有了。”}
}
{
    \eR{[English source not present in the checked-in Bencraft/Joly files.]}
}
{
}

\Verse{259}
{
    \zR{宝玉睡眼蒙？}
}
{
    \eR{[English source not present in the checked-in Bencraft/Joly files.]}
}
{
}

\Verse{260}
{
    \zR{，接在手里也没瞧， 便往地下一撂，道：“你们又来哄我了。”}
}
{
    \eR{[English source not present in the checked-in Bencraft/Joly files.]}
}
{
}

\Verse{261}
{
    \zR{说着只是冷笑。}
}
{
    \eR{[English source not present in the checked-in Bencraft/Joly files.]}
}
{
}

\Verse{262}
{
    \zR{凤姐连忙拾起来道：“这 也就奇了，怎么你没瞧就知道呢？”}
}
{
    \eR{[English source not present in the checked-in Bencraft/Joly files.]}
}
{
}

\Verse{263}
{
    \zR{宝玉也不答言，只管笑。}
}
{
    \eR{[English source not present in the checked-in Bencraft/Joly files.]}
}
{
}

\Verse{264}
{
    \zR{王夫人也进屋里来了， 见他这样，便道：“这不用说了。}
}
{
    \eR{[English source not present in the checked-in Bencraft/Joly files.]}
}
{
}

\Verse{265}
{
    \zR{他那玉原是胎里带来的一宗古怪东西，自然他有 道理。}
}
{
    \eR{[English source not present in the checked-in Bencraft/Joly files.]}
}
{
}

\Verse{266}
{
    \zR{想来这个必是人家见了帖儿，照样儿做的。”}
}
{
    \eR{[English source not present in the checked-in Bencraft/Joly files.]}
}
{
}

\Verse{267}
{
    \zR{大家此时恍然大悟。}
}
{
    \eR{[English source not present in the checked-in Bencraft/Joly files.]}
}
{
}

\Verse{268}
{
    \zR{贾琏在外间屋里听见这个话，便说道：“既不是，快拿来给我问问他去。}
}
{
    \eR{[English source not present in the checked-in Bencraft/Joly files.]}
}
{
}

\Verse{269}
{
    \zR{人家 这样事，他还敢来鬼混！”}
}
{
    \eR{[English source not present in the checked-in Bencraft/Joly files.]}
}
{
}

\Verse{270}
{
    \zR{贾母喝住道：“琏儿，拿了去给他，叫他去罢。}
}
{
    \eR{[English source not present in the checked-in Bencraft/Joly files.]}
}
{
}

\Verse{271}
{
    \zR{那也是 穷极了的人，没法儿了，所以见我们家有这样事，他就想着赚几个钱，也是有的。}
}
{
    \eR{[English source not present in the checked-in Bencraft/Joly files.]}
}
{
}

\Verse{272}
{
    \zR{如今白白的花了钱弄了这个东西，又叫咱们认出来了。}
}
{
    \eR{[English source not present in the checked-in Bencraft/Joly files.]}
}
{
}

\Verse{273}
{
    \zR{依着我倒别难为他，把这块 玉还他，说不是我们的，赏给他几两银子，外头的人知道了，才肯有信儿就送来呢。}
}
{
    \eR{[English source not present in the checked-in Bencraft/Joly files.]}
}
{
}

\Verse{274}
{
    \zR{要是难为了这一个人，就有真的人家也不敢拿了来了。”}
}
{
    \eR{[English source not present in the checked-in Bencraft/Joly files.]}
}
{
}

\Verse{275}
{
    \zR{贾琏答应出去。}
}
{
    \eR{[English source not present in the checked-in Bencraft/Joly files.]}
}
{
}

\Verse{276}
{
    \zR{那人还等 着呢，半日不见人来，正在那里心里发虚，只见贾琏气忿忿走出来了。}
}
{
    \eR{[English source not present in the checked-in Bencraft/Joly files.]}
}
{
}

\Verse{277}
{
    \zR{未知如何，下回分解。}
}
{
    \eR{[English source not present in the checked-in Bencraft/Joly files.]}
}
{
}

\Verse{278}
{
    \zR{(本章完)}
}
{
    \eR{[English source not present in the checked-in Bencraft/Joly files.]}
}
{
}
